\documentclass[UTF8]{ctexart}
\usepackage{geometry}
\usepackage{listings}
\usepackage{xcolor}
\usepackage{graphicx}
\usepackage{amsmath}
\usepackage{float}
\usepackage{hyperref}
\usepackage{tcolorbox}

\geometry{a4paper, left=2.5cm, right=2.5cm, top=2.5cm, bottom=2.5cm}

\title{\textbf{A Simple Calculator}}
\author{
    Student ID: 12312710 \\
    Name: 苍熠/Brighton \\
    Course: CS219: Advanced Programming \\
    Date: \today
}
\date{}

\lstset{
    basicstyle=\ttfamily\small,
    keywordstyle=\color{blue},
    commentstyle=\color{gray},
    stringstyle=\color{red},
    frame=single,
    breaklines=true,
    numbers=left,
    numberstyle=\tiny\color{gray},
    language=C
}

\begin{document}

\maketitle
\tableofcontents
\newpage

\section{前言}
你知道的，于C++的5个Project一直是南科大大佬们秀技的修罗场。在GitHub上看几位学长的report，说实话读完之后冥冥中产生了一种敬畏感。但转机恰巧在这个学期出现了，“本来二十几天前就该写完的Project，硬是让DeepSeek搞成样子”。显然，这一次project的侧重方向悄然的发生了某些改变，我们如何使用AI辅助coding成为了我们必须掌握的一种技能，而这次project会不可避免的考查我们这个部分的能力。因此，本project在完成基本要求和部分进阶要求的前提下，融入了大量的科学使用AI工具能力的内容。

\section{功能展示}

\subsection{运行环境}
本程序是由 C 语言编写。在 Linux (Ubuntu / WSL) 系统中，使用 gcc 编译器进行了编译和运行。
\begin{verbatim}
$ gcc main.c -o main
$ ./main
\end{verbatim}

\subsection{交互式输入输出}
程序支持交互式命令行 (REPL) 模式，用户可以直接输入表达式，程序会即时输出结果。支持 `quit` 或 `exit` 命令退出。
\begin{verbatim}
> 1 + 1
1 + 1 = 2
> 100 / 3
100 / 3 = 33.333333333333333333
\end{verbatim}

\subsection{支持多种数字输入形式}
输入解析器非常鲁棒，支持以下格式：
\begin{itemize}
    \item 普通整数：`123`
    \item 小数：`123.456`
    \item 科学计数法：`1e10`, `1.23e-5`, `1E+100` (甚至支持极大指数如 `1e100`)
    \item 带有冗余空格的输入：`  1   +   2  `
\end{itemize}

\subsection{对无限精度的数字进行无精度损耗的计算}
本计算器核心采用动态数组实现的 `NumericVar` 结构，理论上精度仅受限于内存大小。无论是 `1e100 * 1e100` 产生的大数，还是高精度除法，都能保持无损或用户指定精度的计算。

\subsection{支持部分函数调用}
内置了开根、乘方函数，支持直接在计算中使用。调用方式为 `sqrt(x)`, `x^n`。注意 `sqrt` 默认保留 18 位小数。
\begin{tcolorbox}[colback=gray!10, colframe=gray!50, arc=5pt]
\begin{verbatim}
> sqrt(10)
sqrt(10) = 3.162277660168379332
> 2 ^ 10
2 ^ 10 = 1024
> 3 ^ 100 
3 ^ 100 = 515377520732011331036461129765621272702107522001
\end{verbatim}
\end{tcolorbox}

\subsection{对无限精度的数字进行任意精度调整和计算}
对于除法和开方运算，由于可能产生无限循环小数，计算器采用了**固定高精度策略**。目前默认保留小数点后 **18位** 有效数字，这足以满足绝大多数科学计算需求。对于加减乘法，则保持**完全精度**，不会发生溢出或精度丢失。
\begin{tcolorbox}[colback=gray!10, colframe=gray!50, arc=5pt]
\begin{verbatim}
> 1 / 3
1 / 3 = 0.333333333333333333
> sqrt(2)
sqrt(2) = 1.414213562373095048
> 123456789123456789 * 123456789123456789
15241578780673678515622620750190521
\end{verbatim}
\end{tcolorbox}

\subsection{错误输入识别}
对于非法输入（如除以零、格式错误的数字、不支持的运算符），程序会捕捉并打印友好的错误信息，而不会导致程序崩溃。
\begin{verbatim}
> 1 / 0
Error: 除数不能为零。
> 1 + abc
Error: 输入无法解释为数字。
\end{verbatim}

\section{功能实现}

\subsection{高精度}
\subsubsection{分析}
为了支持任意精度的浮点运算，仅仅使用 C 语言内置的 `double` 或 `long double` 是不够的。我们需要自己管理数字的每一位。

\subsubsection{数据存储方式}
采用了类似 Python `long` 对象的实现方式，使用一个数组存储数字的“段”（digits），每一段代表 1000 进制（BASE=1000）的一位。
\begin{lstlisting}
typedef struct { 
    unsigned int *digits; // 动态数组，存储 BASE 进制的位
    int ndigits;          // 数组长度
    int weight;           // 权值：小数点相对于 digits[0] 的位置
    int sign;             // 符号：1 正, -1 负, 0 零
    int scale;            // 显示精度
} NumericVar;
\end{lstlisting}

一个封装的数字 `NumericVar` 共有以下关键信息：存储每一位的动态数组 \texttt{digits[]}，一个整数变量存储数组的长度 \texttt{ndigits}，一个整数变量存储相对于小数点的权重 \texttt{weight}，以及符号标志 \texttt{sign}。为了计算方便，我们把最高有效位 (MSB) 放在了数组的第一位 (\texttt{digits[0]})。

我们的数值实质上可以表示为：
\[
\text{Value} = \text{sign} \times \left( \sum_{i=0}^{\text{ndigits}-1} \text{digits}[i] \times \text{BASE}^{\text{weight}-i} \right)
\]
其中 \(\text{BASE} = 1000\)，\(\text{digits}[i] \in \{0, 1, \dots, \text{BASE}-1\}\)。
这种设计也是对于浮点数 `Sign-Significand-Exponent` 模型的变体。

\subsection{加减运算}
\subsubsection{原理}
加减法的核心在于**对齐小数点**。由于我们使用 `weight` 来表示位权，两个数相加前，必须先把它们的指数（weight）对齐，通常是将精度较低的数对齐到精度较高的数（即补零）。
\subsubsection{实现}
\begin{lstlisting}[language=C, title=Add(A, B)]
// 伪代码逻辑: Add(A, B)
1. 如果 A.符号 != B.符号:
    // 异号相加转为减法
    如果 |A| > |B| 返回 SubAbs(A, B)，符号为 A.符号
    否则 返回 SubAbs(B, A)，符号为 B.符号

2. // 同号相加 (AddAbs)
   初始化 R，分配足够空间
   设 R.权重 = max(A.权重, B.权重) + 1  // 预留进位空间
   进位 carry = 0

3. // 从低位向高位遍历
   对于 i 从 0 到 (max_weight - min_weight):
       val_a = A 在当前权重位的数值 (若无则为0)
       val_b = B 在当前权重位的数值 (若无则为0)
       
       sum = val_a + val_b + carry
       
       R.digits[i] = sum % BASE  // 当前位
       carry = sum / BASE        // 计算进位

4. 去除 R 的前导零和无效零
5. 返回 R
\end{lstlisting}

\begin{lstlisting}[language=C, title=SubAbs(A, B)]
// 伪代码逻辑: SubAbs(A, B) (前置条件 |A| >= |B|)
1. 初始化 R，R.权重 = A.权重
   借位 borrow = 0

2. // 从低位向高位遍历
   对于 i 从 0 到 R.长度:
       val_a = A 在当前权重位的数值
       val_b = B 在当前权重位的数值
       
       diff = val_a - val_b - borrow
       
       如果 diff < 0:
           diff = diff + BASE  // 向高位借位
           borrow = 1
       否则:
           borrow = 0
           
       R.digits[i] = diff

3. 去除 R 的前导零，返回 R
\end{lstlisting}


\subsection{乘法算法}
\subsubsection{原理}
实现了经典的长乘法（Long Multiplication），时间复杂度为 $O(N \times M)$。虽然不是最快的 Karatsuba 或 FFT，但对于常规计算器应用足够高效且易于实现。
\subsubsection{实现}
\begin{lstlisting}[language=C, title=Multiply(A, B)]
// 伪代码逻辑: Multiply(A, B)
1. 初始化 R，分配空间 = A.长度 + B.长度
   R.digits 全部置零
   R.权重 = A.权重 + B.权重 + 1

2. // 长乘法 (Long Multiplication)
   对于 i 从 0 到 A.长度 - 1:
       如果 A.digits[i] == 0: continue
       
       进位 carry = 0
       对于 j 从 0 到 B.长度 - 1:
           // 计算在结果数组中的位置
           k = i + j + 1
           
           // 累加乘积：当前位 = 旧值 + (A位 * B位) + 进位
           sum = R.digits[k] + (A.digits[i] * B.digits[j]) + carry
           
           R.digits[k] = sum % BASE
           carry = sum / BASE
       
       // 处理当前行剩余的进位
       R.digits[i] += carry

3. R.符号 = A.符号 * B.符号
4. 去除 R 的前导零，返回 R
\end{lstlisting}
特别注意了 `1e100` 这种科学计数法的大数乘法，通过调整 `weight` 而不是移动内存来优化性能。

\subsection{除法算法}
\subsubsection{原理}
实现了 Knuth 的 **Algorithm D** (The Art of Computer Programming, Vol 2)，这是一种能够处理多精度整数除法的标准算法。
\subsubsection{代码实现}
\begin{lstlisting}[language=C, title=Divide(P / Q)]
// 伪代码逻辑: Divide(A, B) (Knuth Algorithm D)
1. 归一化: 确保 B 的最高位 >= BASE/2 (提高估商精度)
   Q.权重 = A.权重 - B.权重
   dividend = A (作为余数缓冲区)

2. // 主循环: 从高位向低位试商
   对于 shift 从 Q.权重 递减到 0:
       // 估算商 digit (q_hat)
       // 使用 dividend 的最高2位 (u0, u1) 和 B 的最高1位 (v1)
       q_hat = (u0 * BASE + u1) / v1
       如果 q_hat >= BASE: q_hat = BASE - 1
       
       // 修正 q_hat：如果 (B * q_hat) > 当前对应的 dividend 片段
       当 (B * q_hat) > current_dividend_window:
           q_hat = q_hat - 1
       
       // 记录商
       Q.digits[对应下标] = q_hat
       
       // 乘减操作：更新 dividend (余数)
       dividend = dividend - (B * q_hat * BASE^shift)

3. 返回 Q, dividend (此时即为 Rem)
\end{lstlisting}
过程包含：归一化、估商、修正、去归一化。

\subsection{开方算法}
\subsubsection{原理}
采用 **以整数为基础的牛顿迭代法 (Newton-Raphson Method)**。
迭代公式：
$$ x_{n+1} = \frac{1}{2} (x_n + \frac{a}{x_n}) $$
为了处理小数，先将数字放大 $10^{2k}$ 倍变成整数开方，最后再缩小。

\begin{lstlisting}[language=C, title=Sqrt(N)]
// 伪代码逻辑: Sqrt(N, scale)
1. // 放大转换
   M = N * 10^(2 * scale)
   
2. // 初始猜测 X_0
   根据 M 的十进制位数估算，X_0 = 10^(digits(M) / 2)

3. // 牛顿迭代循环
   重复直到收敛 (X_new >= X_old):
       X_new = (X_old + M / X_old) / 2
       更新 X_old = X_new

4. // 结果修正 (确保是 floor(sqrt(M)))
   如果 (X+1)^2 <= M: X = X + 1
   如果 X^2 > M: X = X - 1

5. // 缩小还原
   结果 = X * 10^(-scale)
   返回 结果
\end{lstlisting}
牛顿法具有二次收敛速度，对于高精度计算非常有效。

\subsection{指数计算}
\begin{lstlisting}[language=C, title=Power(A, B)]
// 伪代码逻辑: Power(Base, Exp)
1. R = 1
2. 循环 i 从 1 到 Exp:
       R = Multiply(R, Base) // 调用长乘法
3. 返回 R
\end{lstlisting}
（注：当前实现为线性循环，未来可优化为 $O(\log N)$ 的快速幂算法）。

\section{我是怎么使用 AI 的}

\subsection{引言}
在本项目中，GitHub Copilot 扮演了结对编程伙伴的角色。

\subsection{Copilot 也是神}
\subsubsection{自动补全}
在编写 `NumericVar` 结构体和基础的内存管理函数（如 `alloc_var`）时，Copilot 能够根据上下文极其准确地补全代码，几乎不需要修改。

\subsubsection{Fix Bug}
在测试阶段，我发现 `1e100 * 1e100` 的结果不正确（算出了 100）。我直接将问题描述给 Copilot Agent，它迅速定位到了 `mul_var_pow10_inplace` 函数中对于 `digits` 数组移动逻辑的冗余和错误，并给出了精准的修复方案（只修改 weight，不移动内存）。

\subsubsection{Test Case 生成}
为了验证计算器的鲁棒性，我让 Copilot 生成了一个 `run_500_tests.py` 脚本。这个脚本利用 Python 的 `decimal` 模块作为“标准答案”，随机生成了 500 组涵盖整数、小数、科学计数法、边缘情况（如除以零）的测试用例，极大地提升了测试效率。

\section{不足之处}
\begin{itemize}
    \item **乘方功能限制**：目前 `^` 运算符仅支持整数指数，不支持浮点数指数（如 `2^0.5`）。
    \item **函数库不够丰富**：尚未实现 `sin`, `cos`, `log` 等三角函数和对数函数。
    \item **性能优化**：对于超大规模数字（如几万位）的乘法，O(N^2) 的算法由于没有实现 FFT 会变慢。
\end{itemize}

\section{结语}
通过本项目，我不仅深入理解了高精度运算的底层原理（位权管理、进位处理、数值分析算法），更学会了如何高效地与 AI 协作。从代码生成到 Bug 修复，再到自动化测试体系的构建，AI 工具极大地提升了工程效率，让我能将更多精力集中在核心算法逻辑的设计上。

\begin{thebibliography}{99}
    \bibitem{knuth} Knuth, D. E. (1998). \textit{The Art of Computer Programming, Volume 2: Seminumerical Algorithms} (3rd ed.). Addison-Wesley.
    \bibitem{newton} Ypma, T. J. (1995). Historical development of the Newton-Raphson method. \textit{SIAM Review}, 37(4), 531-551.
    \bibitem{postgres} The PostgreSQL Global Development Group. (n.d.). PostgreSQL Source Code (numeric.c). Retrieved from https://github.com/postgres/postgres
    \bibitem{copilot} GitHub. (n.d.). GitHub Copilot - Your AI pair programmer. Retrieved from https://github.com/features/copilot
\end{thebibliography}

\end{document}
